% ==============================================================================
%  INFORME - Análisis comparativo BEST (Y=1) vs WORST (Y=0)
%  Dataset: d9_strong.txt dividido por la variable Y
%  Formato: APA 7ma edición (manual, tipo informe de tesis)
%  Compilar en Overleaf con pdflatex (2-3 pasadas)
% ==============================================================================

\documentclass[12pt, a4paper]{article}

% --- PAQUETES ---
\usepackage[utf8]{inputenc}
\usepackage[T1]{fontenc}
\usepackage[spanish, es-tabla]{babel}
\usepackage{amsmath, amssymb}
\usepackage{graphicx}
\usepackage{booktabs}
\usepackage{hyperref}
\usepackage{url}
\usepackage{float}
\usepackage{algorithm}
\usepackage{algpseudocode}
\usepackage{array}
\usepackage{setspace}
\usepackage{indentfirst}
\usepackage{fancyhdr}
\usepackage{titlesec}
\usepackage{titling}
\usepackage{geometry}
\usepackage{times}
\usepackage{caption}
\usepackage{tocloft}

% --- FORMATO APA 7 ---
\geometry{margin=2.54cm}
\doublespacing
\setlength{\parindent}{1.27cm}
\setlength{\parskip}{0pt}

% Encabezado (running head)
\pagestyle{fancy}
\fancyhf{}
\fancyhead[R]{\thepage}
\renewcommand{\headrulewidth}{0pt}

% Títulos de sección (APA 7)
\titleformat{\section}[block]{\raggedright\bfseries\fontsize{12}{14}\selectfont}{\thesection}{1em}{}
\titlespacing{\section}{0pt}{12pt}{6pt}

\titleformat{\subsection}[block]{\bfseries\fontsize{12}{14}\selectfont\raggedright}{\thesubsection}{1em}{}
\titlespacing{\subsection}{0pt}{12pt}{6pt}

% Captions estilo APA 7
\captionsetup{font={small,sf}, labelfont={bf}, textfont={it}, labelsep=period}

% --- DATOS DEL DOCUMENTO ---
\title{Análisis Comparativo de la Topología del Árbol de Expansión entre
       los Grupos BEST y WORST usando Información Mutua y Algoritmos de
       Prim-Kruskal}
\author{}
\date{}

\begin{document}

% ==============================================================================
% PORTADA (APA 7 manual)
% ==============================================================================

% ==============================================================================
% CARATULA - UNIVERSIDAD NACIONAL DEL ALTIPLANO
% ==============================================================================

\begin{titlepage}
\centering

{\large \textbf{UNIVERSIDAD NACIONAL DEL ALTIPLANO}}\\[0.3cm]
{\large \textbf{FACULTAD DE INGENIERIA MECANICA ELECTRICA, ELECTRONICA Y SISTEMAS}}\\[0.3cm]
{\large \textbf{ESCUELA PROFESIONAL DE INGENIERIA DE SISTEMAS}}\\[1cm]

\includegraphics[width=3.5cm]{logo.png}\\[0.5cm]

{\Large \textbf{ANALISIS COMPARATIVO DE LA TOPOLOGIA DEL ARBOL DE EXPANSION}}\\[0.3cm]
{\large \textbf{Entre los Grupos BEST y WORST usando Informacion Mutua}}\\[1cm]

\textbf{PRESENTADO POR:}\\
Huanacuni Ccama Jhon Edy\\[0.5cm]

\textbf{CURSO:}\\
Lenguajes de Programacion\\[0.8cm]

\textbf{DOCENTE:}\\
ING. JOSE LUIS JUAREZ RUELAS\\[0.8cm]

\textbf{SEMESTRE ACADEMICO:}\\
2026-I

\end{titlepage}

% ==============================================================================
% INDICE
% ==============================================================================

\newpage
\tableofcontents
\newpage

\section*{Resumen}
\addcontentsline{toc}{section}{Resumen}

Este trabajo presenta un análisis comparativo de la estructura de
dependencias entre variables en dos subconjuntos del dataset
\texttt{d9\_strong.txt}: el grupo BEST (502 registros, $Y = 1$) y el grupo
WORST (498 registros, $Y = 0$). Las variables $x_7$ y $x_8$ resultaron
ser columnas idénticas (0 diferencias en 1000 filas), por lo que se
combinaron en una sola variable $Y$. Para cada grupo se aplicó el
pipeline completo de teoría de la información: entropía de Shannon
$H(X)$ para las 8 variables, Información Mutua $I(X;Y)$ entre los 28
pares, matriz IM simétrica de $7 \times 7$, conversión a distancia
mediante $d = 1 - I/\min(H_i, H_j)$, y Árboles de Expansión Mínima y
Máxima mediante los algoritmos de Prim y Kruskal. Los resultados muestran
que $Y$ se vuelve constante en cada grupo ($H(Y) = 0$), degenerando el
árbol MIN con costo total cero. Sin embargo, el árbol MAX basado en MI
directa revela diferencias estructurales: el costo varía de 2.833722 (BEST)
a 2.809188 (WORST), el hub central cambia de $x_1$ (grado 3, BEST) a $x_9$
(grado 4, WORST), y solo 5 de 7 aristas son comunes entre ambos grupos. La
variable $x_3$ presenta la mayor variación de entropía ($\Delta H = -0.069$)
entre grupos. Estos hallazgos demuestran que las topologías de dependencia
son sensibles al contexto y que las variables de partición no deben incluirse
como features en el análisis de cada subconjunto.

\vspace{0.5cm}
\noindent\textit{Palabras clave:} Entropía de Shannon, Información Mutua,
Árbol de Expansión, Prim, Kruskal, BEST vs WORST, variable de partición.

% ==============================================================================
% ABSTRACT
% ==============================================================================

\section*{Abstract}
\addcontentsline{toc}{section}{Abstract}

This paper presents a comparative analysis of variable dependency structures
between two subsets of the \texttt{d9\_strong.txt} dataset: the BEST group
(502 records, $Y = 1$) and the WORST group (498 records, $Y = 0$). Variables
$x_7$ and $x_8$ were found to be identical columns and were merged into $Y$.
For each group, the complete information-theoretic pipeline was applied:
Shannon entropy $H(X)$, Mutual Information $I(X;Y)$ across 28 pairs, an
$7 \times 7$ IM matrix, distance transformation via $d = 1 - I/\min(H_i, H_j)$,
and Minimum and Maximum Spanning Trees via Prim and Kruskal algorithms. Results
show that $Y$ becomes constant ($H(Y) = 0$), degenerating the MIN tree. The MAX
tree reveals structural differences: cost varies from 2.833722 (BEST) to
2.809188 (WORST), the hub shifts from $x_1$ to $x_9$, and only 6 of 7 edges
are shared. Variable $x_3$ shows the largest entropy variation ($\Delta H =
-0.069$), demonstrating context-sensitive dependency topologies.

\vspace{0.5cm}
\noindent\textit{Keywords:} Shannon entropy, Mutual Information, Spanning Tree,
Prim, Kruskal, BEST vs WORST, partition variable.

\newpage

% ==============================================================================
\section{Introducción}
% ==============================================================================

Este trabajo aborda el análisis del dataset \texttt{d9\_strong.txt}, con
1000 registros y 9 variables categóricas que toman valores discretos en
$\{0, 1, 2, 3\}$.

Un análisis previo detectó que $x_7$ y $x_8$ son columnas idénticas
(0 diferencias en 1000 filas). En consecuencia, se combinaron en una sola
variable $Y = x_7$. A partir de $Y$ se definieron dos grupos: BEST
(502 registros, $Y = 1$) y WORST (498 registros, $Y = 0$).

El objetivo del trabajo es aplicar el pipeline completo de teoría de la
información ---entropía, Información Mutua, matriz IM y Árbol de Expansión
Mínima--- a cada grupo por separado y comparar las topologías de red
resultantes para cuantificar los patrones de cambio en el comportamiento de
las variables según el contexto definido por $Y$.

% ==============================================================================
\section{Objetivos}
% ==============================================================================

El objetivo general del presente trabajo es ejecutar el pipeline completo
de teoría de la información y teoría de grafos sobre el dataset
\texttt{d9\_strong.txt}, particionado en los grupos BEST y WORST, para
identificar y cuantificar los cambios en la topología de dependencias
entre variables según el contexto definido por la variable $Y$.

Para alcanzar este objetivo general, se plantean los siguientes objetivos
específicos:

\begin{enumerate}
    \item Combinar las variables redundantes $x_7$ y $x_8$ en una sola
          variable $Y$, y particionar el dataset en dos grupos:
          BEST ($Y = 1$, 502 registros) y WORST ($Y = 0$, 498 registros).

    \item Calcular la entropía $H(X)$, la Información Mutua
          $I(X;Y)$, la matriz IM de $7 \times 7$ y la matriz de
          distancias para cada grupo por separado.

    \item Construir los Árboles de Expansión Mínima (MIN) y
          Máxima (MAX) mediante los algoritmos de Prim y Kruskal
          para cada grupo.

    \item Comparar las topologías de red resultantes y cuantificar
          las diferencias en: costo total, hub central, aristas
          compartidas y pares con dependencia perfecta.

    \item Identificar las variables cuyo comportamiento cambia
          significativamente entre BEST y WORST, y determinar si la
          variable de partición $Y$ afecta la estructura de
          dependencias.
\end{enumerate}

% ==============================================================================
\section{Marco Teórico}
% ==============================================================================

\subsection{Entropía de Shannon}

La entropía de una variable aleatoria discreta $X$ cuantifica su
incertidumbre promedio y se define como \cite{shannon1948}:

\begin{equation}\label{eq:entropy}
    H(X) = -\sum_{x \in \mathcal{X}} P(x) \cdot \log_2 P(x)
\end{equation}

donde $\mathcal{X}$ es el conjunto de valores que toma $X$. La entropía
máxima se alcanza con una distribución uniforme: $H_{\max} = \log_2(k)$,
donde $k$ es la cantidad de valores distintos de la variable.

\subsection{Información Mutua}

La Información Mutua entre dos variables $X$ e $Y$ mide la cantidad de
información compartida entre ellas \cite{cover2006}:

\begin{equation}\label{eq:mi}
    I(X;Y) = \sum_{x} \sum_{y} P(x,y) \cdot
    \log_2 \left( \frac{P(x,y)}{P(x) \cdot P(y)} \right)
\end{equation}

Siempre se cumple $0 \le I(X;Y) \le \min(H(X), H(Y))$, donde $I = 0$
implica independencia total y $I = \min(H)$ implica dependencia perfecta:
una variable determina completamente a la otra.

\subsection{Conversión a Distancia}

Para construir un árbol de costo mínimo, la medida de similitud IM debe
convertirse en una distancia. Se utiliza la siguiente transformación:

\begin{equation}\label{eq:distance}
    d(X_i, X_j) = 1 -
    \frac{I(X_i; X_j)}{\min\!\big(H(X_i), H(X_j)\big)}
\end{equation}

donde $d \in [0,1]$, con $d = 0$ para dependencia perfecta y $d = 1$
para independencia total. Una propiedad importante de la
Ecuación~\ref{eq:distance} es que cuando una variable es constante
($H = 0$), todas las aristas incidentes tienen $d = 0$, produciendo
un árbol degenerado.

\subsection{Algoritmo de Prim}

El Algoritmo~\ref{alg:prim} construye el MST incrementalmente desde
un nodo inicial, agregando en cada paso la arista de menor peso que
conecte un nodo visitado con uno no visitado \cite{prim1957}.

\begin{algorithm}[H]
\caption{Prim --- Árbol de Expansión Mínima}\label{alg:prim}
\begin{algorithmic}[1]
\Function{Prim}{$D[n][n]$} \Comment{Matriz de distancias}
    \State $V \gets \{0\}$ \Comment{Nodo inicial}
    \State $E \gets \{\,\}$ \Comment{Aristas del árbol}
    \While{$|V| < n$}
        \State $mejor \gets \infty$
        \For{\textbf{cada} $u \in V$}
            \For{\textbf{cada} $v \notin V$}
                \If{$D[u][v] < mejor$}
                    \State $mejor \gets D[u][v]$
                    \State $(u^*, v^*) \gets (u, v)$
                \EndIf
            \EndFor
        \EndFor
        \State $V \gets V \cup \{v^*\}$
        \State $E \gets E \cup \{(u^*, v^*, mejor)\}$
    \EndWhile
    \State \Return $E$
\EndFunction
\end{algorithmic}
\end{algorithm}

\subsection{Algoritmo de Kruskal}

El Algoritmo~\ref{alg:kruskal} ordena todas las aristas y las procesa
utilizando una estructura Union-Find, aceptando aquellas que no forman
ciclos \cite{kruskal1956}.

\begin{algorithm}[H]
\caption{Kruskal con Union-Find}\label{alg:kruskal}
\begin{algorithmic}[1]
\Function{Kruskal}{$D[n][n]$}
    \State $E \gets \{\,\}$
    \State $A \gets$ \Call{Ordenar}{aristas $(i,j)$ por $D[i][j]$ asc.}
    \State $UF \gets$ \Call{UnionFind}{$n$} \Comment{$n$ conjuntos unitarios}
    \For{\textbf{cada} $(d, u, v) \in A$}
        \If{\Call{Find}{$u$} $\neq$ \Call{Find}{$v$}}
            \State \Call{Union}{$u, v$}
            \State $E \gets E \cup \{(u, v, d)\}$
            \If{$|E| = n - 1$} \State \textbf{romper} \EndIf
        \EndIf
    \EndFor
    \State \Return $E$
\EndFunction
\Statex
\Function{Find}{$x$} \Comment{Compresión de ruta}
    \If{$parent[x] \neq x$}
        \State $parent[x] \gets$ \Call{Find}{$parent[x]$}
    \EndIf
    \State \Return $parent[x]$
\EndFunction
\Statex
\Function{Union}{$x, y$}
    \State $rx \gets$ \Call{Find}{$x$}, $ry \gets$ \Call{Find}{$y$}
    \If{$rx = ry$} \State \Return \textbf{falso} \EndIf
    \If{$rank[rx] < rank[ry]$}
        \State $parent[rx] \gets ry$
    \ElsIf{$rank[rx] > rank[ry]$}
        \State $parent[ry] \gets rx$
    \Else
        \State $parent[ry] \gets rx$
        \State $rank[rx] \gets rank[rx] + 1$
    \EndIf
    \State \Return \textbf{verdadero}
\EndFunction
\end{algorithmic}
\end{algorithm}

% ==============================================================================
\section{Metodología}
% ==============================================================================

\subsection{Preparación de los Datos}

\begin{enumerate}
    \item Se carga \texttt{d9\_strong.txt} (1000 $\times$ 9).
    \item Se combinan $x_7$ y $x_8$ (idénticas) en una sola variable
          $Y = x_7$. Las columnas finales son: $x_1$, $x_2$, $x_3$,
          $x_4$, $x_5$, $x_6$, $Y$, $x_9$.
    \item Se separa en dos archivos: \texttt{d9\_strong\_B.csv}
          ($Y = 1$, 502 registros) y \texttt{d9\_strong\_W.csv}
          ($Y = 0$, 498 registros).
\end{enumerate}

\subsection{Cálculo de Entropía}

El Algoritmo~\ref{alg:entropia} implementa la Ecuación~\ref{eq:entropy}.
Recibe una lista de probabilidades y devuelve la entropía en bits.

\begin{algorithm}[H]
\caption{Cálculo de Entropía $H(X)$}\label{alg:entropia}
\begin{algorithmic}[1]
\Function{Entropía}{$probs$}
    \State $total \gets 0$
    \For{\textbf{cada} $p \in probs$}
        \If{$p > 0$}
            \State $total \gets total - p \cdot \log_2(p)$
        \EndIf
    \EndFor
    \State \Return $total$
\EndFunction
\end{algorithmic}
\end{algorithm}

\subsection{Cálculo de Información Mutua}

El Algoritmo~\ref{alg:mi} implementa la Ecuación~\ref{eq:mi}. Construye
la tabla de contingencia contando las ocurrencias conjuntas y acumula
los términos de la suma.

\begin{algorithm}[H]
\caption{Cálculo de Información Mutua $I(X;Y)$}\label{alg:mi}
\begin{algorithmic}[1]
\Function{MI}{$X$, $Y$, $n$}
    \State $jc \gets \{\,\}$ \Comment{Tabla de contingencia}
    \For{$i \gets 1$ \textbf{hasta} $n$}
        \State $jc[(X[i], Y[i])] \gets jc[(X[i], Y[i])] + 1$
    \EndFor
    \State $mi \gets 0$
    \For{\textbf{cada} $(v_x, v_y), c \in jc$}
        \State $p_{xy} \gets c / n$
        \State $p_x \gets$ \Call{Contar}{$X = v_x$} $/\, n$
        \State $p_y \gets$ \Call{Contar}{$Y = v_y$} $/\, n$
        \State $razón \gets p_{xy} / (p_x \cdot p_y)$
        \State $mi \gets mi + p_{xy} \cdot \log_2(razón)$
    \EndFor
    \State \Return $mi$
\EndFunction
\end{algorithmic}
\end{algorithm}

\subsection{Pipeline Completo}

El Algoritmo~\ref{alg:pipeline} resume la secuencia completa de
procesamiento aplicada a cada grupo de forma independiente.

\begin{algorithm}[H]
\caption{Pipeline Completo de Procesamiento}\label{alg:pipeline}
\begin{algorithmic}[1]
\Function{Pipeline}{$datos$}
    \State \Call{Frecuencias}{$datos$} $\to$ \Call{Probabilidades}{}
    \State $H \gets$ \Call{Entropía}{$probs$} \Comment{7 valores}
    \State $MI \gets \{\,\}$ \Comment{21 pares}
    \For{\textbf{cada} par $(X_i, X_j)$, $i < j$}
        \State $MI[i,j] \gets$ \Call{MI}{$X_i$, $X_j$, $n$}
    \EndFor
    \State \Call{MatrizIM}{$MI$} $7 \times 7$ \Comment{Simétrica}
    \State $D \gets$ \Call{Distancia}{$MI$, $H$} \Comment{Ec.~\ref{eq:distance}}
    \State $MST_{min} \gets$ \Call{Prim}{$D$} \Comment{MIN}
    \State $MST_{max} \gets$ \Call{Prim}{$MI$} (criterio máx.) \Comment{MAX}
    \State \Return $H$, $MI$, $D$, $MST_{min}$, $MST_{max}$
\EndFunction
\end{algorithmic}
\end{algorithm}

% ==============================================================================
\section{Resultados}
% ==============================================================================

\subsection{Variable $Y$ como Criterio de Partición}

\textit{Cumplimiento del Objetivo 1:} $Y$ se utilizó exitosamente para
particionar el dataset en BEST ($Y=1$, 502 registros) y WORST ($Y=0$,
498 registros), sin incluirse en los CSVs de análisis (7 variables:
$x_1$--$x_6$, $x_9$).

La Tabla~\ref{tab:entropy} compara las entropías de cada variable entre
ambos grupos. Las diferencias observadas reflejan cambios reales en las
distribuciones según el contexto definido por $Y$.

\begin{table}[H]
    \centering
    \caption{Entropías $H(X)$ en BEST y WORST}
    \label{tab:entropy}
    \small
    \begin{tabular}{cccc}
    \toprule
    Variable & BEST & WORST & $\Delta$ \\
    \midrule
    $x_1$ & 0.7586 & 0.7354 & +0.0232 \\
    $x_2$ & 0.9991 & 0.9986 & +0.0005 \\
    $x_3$ & 0.6741 & 0.7431 & $-$0.0690 \\
    $x_4$ & 0.7475 & 0.6860 & +0.0614 \\
    $x_5$ & 1.4116 & 1.4162 & $-$0.0046 \\
    $x_6$ & 0.9974 & 0.9998 & $-$0.0024 \\
    $x_9$ & 1.7448 & 1.6854 & +0.0593 \\
    \bottomrule
    \end{tabular}
\end{table}

$x_3$ presenta el mayor cambio de entropía ($-0.069$), seguido de $x_4$
($+0.061$) y $x_9$ ($+0.059$), indicando que las distribuciones de estas
variables son las más sensibles al valor de $Y$.

\subsection{Matrices IM y de Distancias}

\textit{} se calcularon exitosamente las
matrices IM y de distancias para ambos grupos, verificando que la fila
y columna $Y$ contienen ceros por ser $Y$ constante.

La Figura~\ref{fig:matrices} muestra los heatmaps de las matrices IM
y de distancias para ambos grupos. Dado que $Y$ es constante, todas las
entradas en la fila y columna $Y$ son cero tanto en la matriz IM como
en la matriz de distancias.

\begin{figure}[H]
    \centering
    \includegraphics[width=\textwidth]{fig_matrices.png}
    \caption{Matrices IM y de Distancias para los grupos BEST y WORST}
    \label{fig:matrices}
\end{figure}

\subsection{Árbol MIN --- Análisis de Distancias}

\textit{Cumplimiento del Objetivo 3:} el árbol MIN se construyó
correctamente mediante Prim y Kruskal para ambos grupos, con costos
reales y comparables (Tabla~\ref{tab:min}).

A diferencia del enfoque anterior que incluía $Y$, al excluir esta variable
de los CSVs el árbol MIN no sufre degeneración. Las aristas con $d=0$
($x_4$--$x_9$ y $x_6$--$x_9$) representan dependencias reales entre
variables, no artefactos de una variable constante.

\begin{table}[H]
    \centering
    \caption{Árboles MIN (distancias) para BEST y WORST}
    \label{tab:min}
    \small
    \begin{tabular}{c@{\qquad}c}
    \toprule
    \textbf{BEST (2.782115)} & \textbf{WORST (2.723224)} \\
    \midrule
    $x_9$--$x_4$ (0.000) & $x_9$--$x_6$ (0.000) \\
    $x_9$--$x_6$ (0.000) & $x_9$--$x_4$ (0.000) \\
    $x_5$--$x_2$ (0.344) & $x_5$--$x_2$ (0.312) \\
    $x_1$--$x_5$ (0.453) & $x_1$--$x_5$ (0.424) \\
    $x_5$--$x_9$ (0.989) & $x_9$--$x_3$ (0.992) \\
    $x_1$--$x_3$ (0.997) & $x_5$--$x_9$ (0.995) \\
    \bottomrule
    \end{tabular}
\end{table}

Este resultado es matemáticamente correcto y muestra que ambos árboles
tienen costos cercanos pero diferentes ($\Delta = 0.059$), reflejando
cambios en las dependencias entre grupos.

\subsection{Árbol MAX --- Estructura con MI Directa}

\textit{Cumplimiento del Objetivo 3:} el árbol MAX completa el objetivo
al ofrecer una perspectiva complementaria usando MI directa como peso.

El árbol MAX utiliza directamente la Información Mutua como peso de cada
arista. La Tabla~\ref{tab:max} presenta
las aristas resultantes para cada grupo. Las Figuras~\ref{fig:arbol_best}
y~\ref{fig:arbol_worst} muestran los árboles individuales.

\begin{table}[H]
    \centering
    \caption{Árboles MAX (MI directa) para BEST y WORST}
    \label{tab:max}
    \small
    \begin{tabular}{c@{\qquad}c}
    \toprule
    \textbf{BEST (2.833722)} & \textbf{WORST (2.809188)} \\
    \midrule
    $x_9$--$x_6$ (0.9974) & $x_9$--$x_6$ (0.9998) \\
    $x_9$--$x_4$ (0.7475) & $x_5$--$x_2$ (0.6869) \\
    $x_5$--$x_2$ (0.6555) & $x_9$--$x_4$ (0.6860) \\
    $x_1$--$x_5$ (0.4150) & $x_1$--$x_5$ (0.4237) \\
    $x_5$--$x_9$ (0.0162) & $x_5$--$x_9$ (0.0069) \\
    $x_1$--$x_3$ (0.0022) & $x_9$--$x_3$ (0.0058) \\
    \bottomrule
    \end{tabular}
\end{table}

\begin{figure}[H]
    \centering
    \includegraphics[width=\textwidth]{fig_arbol_best.png}
    \caption{Árbol BEST: MIN (distancias) + MAX (MI directa)}
    \label{fig:arbol_best}
\end{figure}

\begin{figure}[H]
    \centering
    \includegraphics[width=\textwidth]{fig_arbol_worst.png}
    \caption{Árbol WORST: MIN (distancias) + MAX (MI directa)}
    \label{fig:arbol_worst}
\end{figure}

\subsection{Comparación BEST vs WORST}

\textit{Cumplimiento de los Objetivos 4 y 5:} se cuantificaron las
diferencias topológicas entre ambos grupos (costos MIN distintos,
83\% de aristas compartidas) y se identificó a $x_3$ como la variable más
sensible al contexto de $Y$ ($\Delta H = -0.069$).

La Tabla~\ref{tab:comparacion} resume las métricas comparativas más
relevantes entre ambos grupos. La Figura~\ref{fig:comparacion} muestra
los cuatro árboles de expansión lado a lado.

\begin{table}[H]
    \centering
    \caption{Métricas comparativas: BEST vs WORST}
    \label{tab:comparacion}
    \small
    \begin{tabular}{lcc}
    \toprule
    Métrica & BEST & WORST \\
    \midrule
    Registros & 502 & 498 \\
    Costo Prim MIN & 2.782115 & 2.723224 \\
    Costo Kruskal MIN & 2.782115 & 2.723224 \\
    Costo Prim MAX & 2.833722 & 2.809188 \\
    Costo Kruskal MAX & 2.833722 & 2.809188 \\
    Aristas MIN compartidas & \multicolumn{2}{c}{5/6 (83\%)} \\
    \bottomrule
    \end{tabular}
\end{table}

Se destaca que los costos MIN difieren entre BEST (2.782115) y WORST
(2.723224), reflejando un cambio en las dependencias según el contexto
de $Y$. En el árbol MIN, 5 de las 6 aristas son comunes, con una única
arista diferente entre ambos grupos ($x_1$--$x_3$ en BEST vs $x_9$--$x_3$
en WORST).

\begin{figure}[H]
    \centering
    \includegraphics[width=\textwidth]{fig_comparacion.png}
    \caption{Comparación final: BEST vs WORST, MIN y MAX}
    \label{fig:comparacion}
\end{figure}

\subsection{Selección de Variables: Método rVP}

El método del Rango de Variación Permitido (rVP) evalúa qué tan sensible
es la estructura de dependencias a la partición del dataset. Para cada
variable se calcula la distancia de camino dentro del MST (MAX, MI) hacia
todas las demás variables:

\begin{equation*}
    D_{\text{MST}}(X_i, X_j) = \sum_{e \in \text{camino}(X_i \to X_j)} w(e)
\end{equation*}

\begin{equation*}
    \Sigma(X_i) = \sum_{j} D_{\text{MST}}(X_i, X_j)
\end{equation*}

\begin{equation*}
    \Delta(X_i) = \Sigma_{\text{BEST}}(X_i) - \Sigma_{\text{WORST}}(X_i)
\end{equation*}

Una variable es \textbf{crítica} si $|\Delta(X_i)| > 0.5$, indicando que
la partición alteró significativamente su posición topológica en el árbol.
Este umbral es el utilizado por trabajos similares en el curso.

\subsubsection{Resultados del Ranking rVP}

\begin{table}[H]
    \centering
    \caption{Ranking rVP (Rango de Variación Permitido)}
    \label{tab:seleccion}
    \small
    \begin{tabular}{cccccl}
    \toprule
    Pos. & Variable & $\Sigma$ BEST & $\Sigma$ WORST & $\Delta$ & ¿Crítica? \\
    \midrule
    1 & $x_3$  & 2.8521 & 4.5371 & $-1.6850$ & \textbf{Sí} \\
    2 & $x_4$  & 6.2531 & 7.0345 & $-0.7814$ & \textbf{Sí} \\
    3 & $x_9$  & 2.8230 & 3.2972 & $-0.4742$ & No \\
    4 & $x_6$  & 7.8221 & 8.2844 & $-0.4623$ & No \\
    5 & $x_5$  & 2.8299 & 3.2811 & $-0.4511$ & No \\
    6 & $x_1$  & 4.9485 & 4.5259 & $+0.4225$ & No \\
    7 & $x_2$  & 6.2646 & 6.5584 & $-0.2938$ & No \\
    \bottomrule
    \end{tabular}
\end{table}

$x_3$ es, por mucho, la variable más afectada por la partición
($\Delta = -1.685$, 3.4 veces el umbral). En BEST se sitúa cerca del
centro del árbol (conectada a $x_1$), mientras que en WORST se desplaza
a la periferia conectándose a $x_9$. Esta diferencia se explica porque
la partición por $Y$ segmenta los valores de $x_3$.

\begin{figure}[H]
    \centering
    \includegraphics[width=\textwidth]{fig_rvp.png}
    \caption{Ranking rVP: barras de $|\Delta|$ con umbral de criticidad 0.5}
    \label{fig:rvp}
\end{figure}

$x_4$ es la segunda variable crítica ($\Delta = -0.781$). Aunque mantiene
su conexión directa con $x_9$, el peso del enlace cambia entre grupos
y la estructura general del árbol redistribuye las distancias acumuladas.

\subsubsection{Conclusiones del rVP}

Desde la perspectiva del método rVP, la partición por $Y$ \textbf{no es
inocua}: dos de las siete variables ($x_3$ y $x_4$) superan el umbral de
criticidad ($|\Delta| > 0.5$), evidenciando que el dataset original contiene
subpoblaciones con estructuras de dependencia diferenciadas. Las cinco
variables restantes ($x_1$, $x_2$, $x_5$, $x_6$, $x_9$) muestran robustez
estructural.

\subsection{Método Complementario: Matriz de Distancias}

Como enfoque complementario, se aplicó el método de suma de fila completa
de la matriz de distancias $d = 1 - I/\min(H)$, propuesto por el docente.
Para cada variable $X_i$:

\begin{equation*}
    \Sigma_{\text{BEST}}(X_i) = \sum_{j=1}^{7} d_{\text{BEST}}(X_i, X_j)
\end{equation*}
\begin{equation*}
    \Delta(X_i) = \Sigma_{\text{BEST}}(X_i) - \Sigma_{\text{WORST}}(X_i)
\end{equation*}

\begin{table}[H]
    \centering
    \caption{Ranking por suma de fila completa de matriz de distancias}
    \small
    \begin{tabular}{cccc}
    \toprule
    Pos. & Variable & $\Sigma$ BEST & $\Sigma$ WORST & $\Delta$ \\
    \midrule
    1 & $x_5$  & 4.7664 & 4.7269 & $+0.0395$ \\
    2 & $x_2$  & 5.3285 & 5.2997 & $+0.0288$ \\
    3 & $x_1$  & 5.4289 & 5.4077 & $+0.0212$ \\
    4 & $x_6$  & 4.9778 & 4.9975 & $-0.0197$ \\
    5 & $x_3$  & 5.9942 & 5.9815 & $+0.0127$ \\
    6 & $x_9$  & 3.9739 & 3.9802 & $-0.0063$ \\
    7 & $x_4$  & 4.9896 & 4.9880 & $+0.0016$ \\
    \bottomrule
    \end{tabular}
\end{table}

En este enfoque se presenta el ranking completo sin umbral predefinido,
dejando al lector la interpretación de los resultados. Las variables $x_5$
y $x_2$ encabezan el ranking con los mayores $|\Delta|$, mientras que $x_4$
muestra $|\Delta| \approx 0$.

\begin{figure}[H]
    \centering
    \includegraphics[width=\textwidth]{fig_profesor.png}
    \caption{Ranking del método de matriz de distancias (sin umbral)}
    \label{fig:profesor}
\end{figure}

\subsection{Comparación entre Métodos}

\begin{table}[H]
    \centering
    \caption{Comparación de rankings: rVP vs Matriz de Distancias}
    \small
    \begin{tabular}{cccc}
    \toprule
    Variable & Rank rVP & Rank Distancias & Interpretación \\
    \midrule
    $x_3$  & \textbf{1} (crítica) & 5 & Cambia posición en el árbol, no peso global \\
    $x_4$  & \textbf{2} (crítica) & 7 & Cambia rol topológico, peso estable \\
    $x_5$  & 5 & \textbf{1} & Peso global cambia, posición estable \\
    $x_2$  & 7 & \textbf{2} & Peso global cambia, posición estable \\
    \bottomrule
    \end{tabular}
\end{table}

Ambos métodos son complementarios: rVP identifica variables que cambian
su \textbf{rol topológico} en el árbol, mientras que el método de distancias
identifica variables que cambian su \textbf{peso global} de conexión.

% ==============================================================================
\section{Discusión}
\section{Discusión}
% ==============================================================================

\subsection{Papel de la Variable de Partición $Y$}

La variable $Y$ ($x_7$ = $x_8$, columnas idénticas) se utilizó exclusivamente
para particionar el dataset en BEST ($Y=1$) y WORST ($Y=0$), sin incluirse
como variable de análisis en los CSVs resultantes. Esta decisión
metodológica evita la degeneración del árbol MIN que ocurriría si $Y$
---al ser constante en cada grupo--- se incluyera como feature, produciendo
$d=0$ en todas sus aristas incidentes. Al excluir $Y$, ambos árboles MIN
obtienen costos reales y comparables: 2.782115 en BEST y 2.723224 en WORST.

\subsection{Cambio Topológico BEST vs WORST}

Aunque 5 de las 6 aristas (83\%) son comunes entre ambos árboles MIN,
la diferencia en una arista indica que las relaciones de dependencia se
modifican según el contexto de $Y$. El costo MIN difiere ($\Delta = 0.059$),
confirmando este cambio estructural. En el árbol MAX, la diferencia de
costos es menor ($\Delta = 0.025$) pero también se observan diferencias
en las aristas seleccionadas.

\subsection{Implicaciones Prácticas}

El análisis demuestra los siguientes puntos:

\begin{enumerate}
    \item La eliminación de variables redundantes ($x_7$-$x_8 \to Y$)
          simplifica el modelo sin pérdida de información.
    \item La partición por $Y$ revela patrones de cambio en las
          dependencias, cuantificables a través de los árboles MIN y MAX.
    \item La exclusión de $Y$ como variable de análisis es necesaria
          para evitar la degeneración del árbol MIN, ya que $Y$ se vuelve
          constante en cada grupo al ser la variable de partición.
\end{enumerate}

% ==============================================================================
\section{Conclusiones}
% ==============================================================================

\begin{enumerate}
    \item Se implementó exitosamente el pipeline BEST vs WORST,
          generando dos datasets de 502 y 498 registros con 7
          variables cada uno (x1-x6, x9).

    \item La variable $Y$ ($x_7=x_8$) se utilizó exclusivamente como
          criterio de partición, sin incluirse en los CSVs de análisis.
          Esto evita la degeneración del árbol MIN y permite obtener
          costos reales: 2.782115 (BEST) y 2.723224 (WORST).

    \item El árbol MIN revela diferencias estructurales: BEST tiene
          costo 2.782115 y WORST 2.723224,
          con 5 de 6 aristas en común (83\%).

    \item $x_3$ es la variable con mayor $\Delta H$ entre grupos
          ($-0.069$), siendo la más sensible al contexto de $Y$.

    \item Los algoritmos de Prim y Kruskal coinciden en todos los
          casos, validando la implementación.

    \item El método rVP identificó a $x_3$ ($\Delta=-1.685$) y $x_4$
          ($\Delta=-0.781$) como variables críticas (cambio de rol
          topológico), mientras que el método complementario de matriz
          de distancias identificó a $x_5$ ($\Delta=+0.039$) y $x_2$
          ($\Delta=+0.029$) como las de mayor cambio de peso global.
          Ambos enfoques son complementarios y enriquecen el análisis.
\end{enumerate}

% ==============================================================================
\section{Video de Exposición}
% ==============================================================================

El video de exposición de 10 minutos se encuentra disponible en el
siguiente enlace:

\begin{center}
\url{https://drive.google.com/drive/folders/18hHH95C_-iLrdQ3zPD-JWE0-No6Qac50?usp=drive_link}
\end{center}

% ==============================================================================
\section{Repositorio de Código}
% ==============================================================================

El código fuente completo del proyecto se encuentra disponible en:

\begin{center}
\url{https://github.com/liamcruzticona/ENTROPIA_PARTICION}
\end{center}

Archivos del proyecto en \texttt{proyecto\_y\_split/}:
\texttt{generar\_csvs.py}, \texttt{pipeline\_comparar.py},
\texttt{visual\_split.py}, \texttt{d9\_strong\_B.csv},
\texttt{d9\_strong\_W.csv}, \texttt{resultados\_split.xlsx}.

% ==============================================================================
\section*{Apéndices}
\addcontentsline{toc}{section}{Apéndices}
% ==============================================================================

\subsection{Salida de Consola}

La Figura~\ref{fig:consola} muestra la salida de consola del pipeline
comparativo con todos los resultados numéricos.

\begin{figure}[H]
    \centering
    \includegraphics[width=0.60\textwidth]{fig_consola.png}
    \caption{Salida de consola del pipeline comparativo}
    \label{fig:consola}
\end{figure}

\subsection{Hoja de Cálculo Excel}

La Figura~\ref{fig:excel} muestra la hoja COMPARACION del archivo Excel
generado con los resultados completos.

\begin{figure}[H]
    \centering
    \includegraphics[width=0.60\textwidth]{fig_excel.png}
    \caption{Hoja de comparación del archivo Excel generado}
    \label{fig:excel}
\end{figure}

\section*{Referencias}
\addcontentsline{toc}{section}{Referencias}

\begin{thebibliography}{99}

\bibitem{shannon1948}
Shannon, C.~E. (1948).
Una teoría matemática de la comunicación.
\textit{Bell System Technical Journal}, \textit{27}(3), 379--423.

\bibitem{cover2006}
Cover, T.~M., y Thomas, J.~A. (2006).
\textit{Elementos de teoría de la información} (2.\textsuperscript{a} ed.).
Wiley-Interscience.

\bibitem{prim1957}
Prim, R.~C. (1957).
Redes de conexión más cortas y algunas generalizaciones.
\textit{Bell System Technical Journal}, \textit{36}(6), 1389--1401.

\bibitem{kruskal1956}
Kruskal, J.~B. (1956).
Sobre el subárbol de expansión más corto de un grafo y el problema del
agente viajero.
\textit{Proceedings of the American Mathematical Society},
\textit{7}(1), 48--50.

\end{thebibliography}

% ==============================================================================
\end{document}
